\section*{ID7360: Example 2}
\paragraph{Metadata.}
PiID: 3904975008656 | RTEID: RTE:P28580.5145:T3.1096 | UUID: 11cc1d30-677c-597b-bf48-c9253c10b356

\paragraph{Canonical Statement.}
\#\#\# Example 2 image\_group\{"layout":"carousel","aspect\_ratio":"1:1","query":["James Webb star eight diffraction spikes","JWST 8 spike star image","JWST diffraction spike pattern 8 points star"],"num\_per\_query":2\} - Observed spike count: **8 spikes** (six strong + two faint). citeturn0image4turn0image6 - Cause: a hexagonal segmented mirror (6 edges) + secondary-mirror support vanes (3 supports) generating overlapping spike sets. - Equation: Mirror edges produce 6 spikes; support vanes produce 6 spikes but 4 overlap → visible \textasciitilde{} 8. \textbackslash{}[ \textbackslash{}text\{Mirror edges: \}6,\textbackslash{}quad \textbackslash{}text\{support vanes: \}6,\textbackslash{}quad \textbackslash{}text\{overlap: \}4 \textbackslash{};\textbackslash{};\textbackslash{}Longrightarrow\textbackslash{}; 6 + 6 - 4 = 8 \textbackslash{}]

\paragraph{Canonical Equation.}
\[
\text{Mirror edges: }6,\quad \text{support vanes: }6,\quad \text{overlap: }4 \;\;\Longrightarrow\; 6 + 6 - 4 = 8
\]

\paragraph{Dictionary.}
Relation: Mirror edges: 6, support vanes: 6, overlap: 4 ⟹ 6 + 6 - 4 = 8

\paragraph{Encyclopedia.}
Example 2 is an algebraic definition, written Mirror edges: 6, support vanes: 6, overlap: 4 ⟹ 6 + 6 - 4 = 8. It is an algebraic definition expressing the quantity directly in terms of the variables on its right-hand side. Within the Trawin Zero Point registry it serves as one element of the framework's mathematical narrative.
